\documentclass[10pt,a4paper]{ctexart}

\usepackage[margin=0.62in]{geometry}
\usepackage{fontspec}
\usepackage{xeCJK}
\usepackage{microtype}
\usepackage{xcolor}
\usepackage{enumitem}
\usepackage{hyperref}
\usepackage{titlesec}
\usepackage{setspace}
\usepackage{array}
\usepackage{tabularx}

\IfFontExistsTF{TeX Gyre Pagella}{
  \setmainfont{TeX Gyre Pagella}
}{
  \setmainfont{Times New Roman}
}
\IfFontExistsTF{TeX Gyre Heros}{
  \setsansfont{TeX Gyre Heros}
}{
  \setsansfont{Arial}
}
\IfFontExistsTF{Microsoft YaHei}{
  \setCJKmainfont{Microsoft YaHei}
  \setCJKsansfont{Microsoft YaHei}
}{
  \setCJKmainfont{SimSun}
  \setCJKsansfont{SimHei}
}

\definecolor{accent}{HTML}{203247}
\definecolor{accentsoft}{HTML}{49617D}
\definecolor{gold}{HTML}{B08A57}
\definecolor{ruleline}{HTML}{D7CAB7}

\pagenumbering{gobble}
\setlength{\parindent}{0pt}
\setstretch{1.05}
\hypersetup{
  colorlinks=true,
  urlcolor=accent,
  linkcolor=accent,
  pageanchor=false
}

\titleformat{\section}{\large\bfseries\color{accent}\sffamily}{}{0em}{}[\vspace{-0.32em}\color{ruleline}\titlerule]
\titlespacing*{\section}{0pt}{0.78em}{0.32em}

\setlist[itemize]{leftmargin=1.35em, itemsep=0.1em, topsep=0.08em, parsep=0pt, partopsep=0pt, label=\textcolor{gold}{\raisebox{0.25ex}{\tiny$\bullet$}}}

\newcommand{\resumeSubheading}[4]{
  \begin{tabular*}{\textwidth}{@{}>{\raggedright\arraybackslash}p{0.68\textwidth}@{\extracolsep{\fill}}>{\raggedleft\arraybackslash}p{0.27\textwidth}@{}}
  \textbf{#1} & #2 \\
  {\color{accentsoft}\small #3} & {\color{accentsoft}\small #4}
  \end{tabular*}\\[-0.18em]
}

\newcommand{\pubentry}[2]{
  \item #1\\[-0.08em]
  {\color{accentsoft}\small #2}
}

\newcommand{\headerlink}[2]{\href{#1}{\textcolor{white}{#2}}}

\begin{document}

\begingroup
\setlength{\fboxsep}{12pt}
\noindent\colorbox{accent}{%
\parbox{\dimexpr\textwidth-2\fboxsep\relax}{%
  {\fontsize{25}{27}\selectfont\color{white}\bfseries 钱佳辰} \\
  \vspace{0.18em}
  {\small\color{white}香港城市大学计算机科学硕士生} \\
  {\small\color{white!88}Jiachen Qian} \\
  \vspace{0.2em}
  {\footnotesize\color{gold}可信人工智能 \textbar{} 多模态大模型 \textbar{} 智能体安全 \textbar{} 对抗防御} \\
  \vspace{0.35em}
  {\footnotesize \headerlink{tel:+8613862215800}{+86 138 6221 5800} \enspace
  \textcolor{gold}{|} \enspace
  \headerlink{mailto:72510756@cityu-dg.edu.cn}{72510756@cityu-dg.edu.cn} \enspace
  \textcolor{gold}{|} \enspace
  \headerlink{https://qianjiachen.github.io/homepage/}{Homepage} \enspace
  \textcolor{gold}{|} \enspace
  \headerlink{https://scholar.google.com/citations?user=zN0XzvAAAAAJ\&hl=zh-CN}{Google Scholar} \enspace
  \textcolor{gold}{|} \enspace
  \headerlink{https://github.com/qianjiachen}{GitHub}} \\
  \vspace{0.2em}
  \textcolor{gold}{\rule{\linewidth}{0.8pt}}
}}
\endgroup

\vspace{0.5em}

\section*{个人简介}
香港城市大学计算机科学硕士生，本科毕业于南京信息工程大学软件工程专业。研究聚焦于可信多模态人工智能与智能体安全，重点关注购物智能体中的视觉对抗失效、Agentic Recommender 中的长时记忆投毒，以及 LLM 的结构敏感行为分析。已围绕多模态智能体安全完成两篇独立一作论文，参与 1 篇关于大模型顺序敏感性与确定性结构重建的 arXiv 预印本，并有 1 篇面向社交媒体篡改图像识别大模型攻击的 ECCV 在投论文，具备从问题定义、方法设计、实验构建到论文撰写的完整科研训练。

\section*{研究方向}
可信多模态人工智能、多模态大模型智能体、智能体安全、对抗攻击与防御、长时记忆投毒与防护、LLM 行为分析与结构探测、篡改图像识别大模型安全、可靠决策机制。

\section*{教育背景}
\resumeSubheading{香港城市大学}{中国香港}{计算机科学硕士}{2025 年 9 月 -- 2027 年 6 月（预计）}
\begin{itemize}
  \item GPA：3.3/4.0
  \item 当前研究围绕可信多模态智能体、记忆增强 AI 系统以及对抗鲁棒性展开。
\end{itemize}

\resumeSubheading{南京信息工程大学}{南京，中国}{软件工程学士}{2021 年 9 月 -- 2025 年 6 月}

\section*{科研项目}
\resumeSubheading{多模态购物智能体安全研究}{中国香港}{香港城市大学研究生科研项目}{2025 年 9 月 -- 至今}
\begin{itemize}
  \item 研究价格约束场景下多模态购物智能体的可信决策问题，关注不可感知的视觉线索如何覆盖显式文本价格证据。
  \item 在 \textit{Penny Wise, Pixel Foolish} 中提出 Visual Dominance Hallucination 失效现象，并设计 PriceBlind 隐形攻击框架，在保持像素级视觉保真的同时诱导模型违反“购买最便宜商品”的约束。
  \item 在 E-ShopBench 上实现约 80\% 白盒攻击成功率，并在 GPT-4o、Gemini-1.5-Pro、Claude-3.5-Sonnet 上取得 35--41\% 迁移攻击成功率，同时分析 Verify-then-Act 与鲁棒编码器防御效果。
\end{itemize}

\resumeSubheading{长时记忆推荐智能体安全研究}{中国香港}{香港城市大学研究生科研项目}{2025 年 9 月 -- 至今}
\begin{itemize}
  \item 研究 Agentic Recommender 长时规划中的安全风险，关注多模态长时记忆如何成为后续推荐操控的攻击面。
  \item 在 \textit{Visual Inception} 中设计视觉记忆投毒攻击，将休眠触发器植入用户上传图像，并在未来推荐规划阶段劫持推理链而无需 prompt injection。
  \item 提出 CognitiveGuard 双进程防御，结合扩散式输入净化与反事实一致性验证，将 Goal-Hit Rate 从约 85\% 降至约 10\%，并支持不同延迟开销配置。
\end{itemize}

\resumeSubheading{社交媒体篡改图像识别大模型攻击研究}{中国深圳}{深圳北理工莫斯科大学在研项目}{2026 年}
\begin{itemize}
  \item 围绕社交媒体场景中的篡改图像识别大模型开展安全研究，关注识别导向大模型在真实在线传播环境中的对抗脆弱性。
  \item 设计面向篡改图像识别大模型的攻击流程，分析攻击对模型识别可靠性与判别行为的影响。
  \item 已形成 1 篇 ECCV 在投论文，目前处于审稿阶段。
\end{itemize}

\section*{学术亮点}
\begin{itemize}
  \item 已形成两篇独立一作多模态智能体安全论文、1 篇 arXiv 合作预印本与 1 篇 ECCV 在投论文。
  \item 围绕视觉对抗攻击、多模态记忆投毒与智能体可靠决策建立了清晰研究主线。
  \item 能够将安全威胁模型转化为可量化的攻击与防御框架，并在前沿多模态模型与推荐智能体上开展系统评测。
  \item 具备从问题定义、方法设计、实验构建、结果分析到论文撰写的独立科研能力。
\end{itemize}

\section*{论文与稿件}
\textbf{独立一作论文}
\begin{itemize}
  \pubentry{\textbf{Qian, Jiachen} and Zhaolu Kang. ``Penny Wise, Pixel Foolish: Bypassing Price Constraints in Multimodal Agents via Visual Adversarial Perturbations.'' \textit{Findings of ACL 2026}，已录用；\href{https://arxiv.org/abs/2604.16515}{arXiv:2604.16515}。}{独立一作；多模态购物智能体价格约束绕过攻击与防御}
  \pubentry{\textbf{Qian, Jiachen}. ``Visual Inception: Compromising Long-term Planning in Agentic Recommenders via Multimodal Memory Poisoning.'' \textit{ACL 2026 Main Conference}，已录用；\href{https://arxiv.org/abs/2604.16966}{arXiv:2604.16966}。}{独立一作；面向 Agentic Recommender 的长时记忆投毒与防御}
\end{itemize}

\textbf{合作预印本}
\begin{itemize}
  \pubentry{He, Yingjie, Zhaolu Kang, Kehan Jiang, \textbf{Qian, Jiachen}, et al. ``How Order-Sensitive Are LLMs? OrderProbe for Deterministic Structural Reconstruction.'' \textit{arXiv preprint} \href{https://arxiv.org/abs/2601.08626}{arXiv:2601.08626}, 2026.}{共同作者；LLM 顺序敏感性与确定性结构重建}
\end{itemize}

\textbf{在投论文}
\begin{itemize}
  \pubentry{Manuscript on attacks against large models for tampered social-media image detection. Submitted to \textit{ECCV}, under review, 2026.}{合作研究；社交媒体篡改图像识别大模型攻击}
\end{itemize}

\section*{学术服务}
\begin{itemize}
  \item 担任 \textit{IEEE Transactions on Information Forensics and Security (TIFS)} 审稿人。
\end{itemize}

\section*{荣誉奖项}
\begin{itemize}
  \item 2025 届：南京信息工程大学本科学院优秀毕业生、校级优秀毕业生
  \item 2021 -- 2022 学年：一等奖学金
  \item 2022 -- 2023 学年：三等奖学金
  \item 2023 -- 2024 学年：三等奖学金
\end{itemize}

\section*{技术能力}
\begin{itemize}
  \item 编程语言：Python、C++、SQL
  \item AI 与科研工具：PyTorch、CUDA、Scikit-learn、LLM/VLM 开发、多模态智能体评测、篡改图像识别模型攻击分析
  \item 科研工程能力：Codex、AI 辅助编程工作流、实验自动化、对抗攻击与防御评测、可复现实验分析
  \item 研究主题：可信人工智能、多模态智能体安全、记忆增强智能体安全、篡改图像识别大模型攻击、对抗攻击与防御
\end{itemize}

\end{document}
